
% ════════════════════════════════════════════
% 参考文献（10.参考文献.tex）—— 写作提示词见本文件末尾【本章节写作提示词】。
% ════════════════════════════════════════════
\newpage
\begin{thebibliography}{99}

\bibitem{jiang2011shuxue} 姜启源, 谢金星, 叶俊. 数学模型[M]. 4版. 北京: 高等教育出版社, 2011.
\bibitem{si2015suanfa} 司守奎, 孙玺菁. 数学建模算法与应用[M]. 北京: 国防工业出版社, 2015.
\bibitem{sheng2008gailv} 盛骤, 谢式千, 潘承毅. 概率论与数理统计[M]. 4版. 北京: 高等教育出版社, 2008.
\bibitem{zhou2016jiqi} 周志华. 机器学习[M]. 北京: 清华大学出版社, 2016.
\bibitem{zhang2004juzhen} 张贤达. 矩阵分析与应用[M]. 北京: 清华大学出版社, 2004.
\bibitem{zhuo2014matlab} 卓金武, 李必文, 魏永生. MATLAB在数学建模中的应用[M]. 2版. 北京: 北京航空航天大学出版社, 2014.
\bibitem{liu2014suiji} 刘次华. 随机过程[M]. 5版. 武汉: 华中科技大学出版社, 2014.
\bibitem{xie2009gps} 谢钢. GPS原理与接收机设计[M]. 北京: 电子工业出版社, 2009.
\bibitem{wang2007cexiang} 王国玉, 王连东. 多站测向交叉定位及其精度分析[J]. 舰船电子对抗, 2007, 30(3): 36-39.
\bibitem{liu2008cexiang} 刘刚, 郭福成. 测向交叉定位系统误差分析[J]. 电子对抗, 2008(1): 21-25.
\bibitem{yang2009dingwei} 杨剑锋, 张贤达. 无线传感器网络节点定位算法研究进展[J]. 计算机学报, 2009, 32(9): 1710-1725.

\end{thebibliography}
